% ============================================================================
% Cross-country synthesis (T-R8) and scenario discovery (PRIM).
% Pre-registered 2026-10-14 (OQ-050 closure; OQ-056), before any
% real-country result was produced.
% Notation follows docs/METHODOLOGY.md and the H1 to H5 sections.
% Citation keys marked "verify" must be checked before submission.
% ============================================================================

\section{Cross-Country Synthesis of the Hypotheses}
\label{sec:synthesis}

\subsection{Purpose}
\label{subsec:synthesis_purpose}

The five hypotheses produce verdicts at different levels. H1, H2, H3 and H5 are
evaluated for each country and technology, giving six verdicts each for Brazil,
Portugal and India with solar photovoltaics and onshore wind. H4 is comparative
and gives one verdict per technology. The synthesis answers the question of
generalisation (RQ5): which hypotheses are observed in which national contexts,
and whether the pattern is the same across them.

The synthesis reports what was observed; it does not compress the verdicts into
a single score. With three countries and two technologies, a majority rule would
turn a small number of cases into an apparently general conclusion, and the
choice of its threshold would weigh more than the evidence. The counts and the
consistency of the verdicts are therefore reported as they are.

\subsection{Country Profiles}
\label{subsec:synthesis_profiles}

For each country and technology, the profile gathers the verdict of each
hypothesis and the numbers on which it rests:
\begin{itemize}
    \item H1: $C_k$ and $V_k$ with the reference value $p_k^2$, and the verdict;
    \item H2: $O$, $O_{\mathrm{ref}}$ and the robustness penalty $\Delta_z$;
    \item H3: $D$, $D_{\mathrm{ref}}$, and the price and benefit of robustness
    $(\Pi,B)$;
    \item H4: $D_{\mathrm{clim}}$, $D_{\mathrm{tec}}$, $\rho$ and the country
    profile, with the technology-level verdict;
    \item H5: $E_0$, $G$, $H_0$, $R_{\mathrm{loss}}$ and $R_{\mathrm{exp}}$;
\end{itemize}
together with the size of the candidate set, the shortlist size $k$ and, for
every number that has one, its range across the sensitivity variants
(Section~\ref{subsec:synthesis_ranges}). The profile is the unit from which the
thesis chapters, the articles derived from them and, later, a per-country page
of an interactive platform are built; it is produced once, in a structured form,
from the persisted results.

\subsection{Cross-Country Summary}
\label{subsec:synthesis_summary}

For each hypothesis and each technology, the summary reports the number of
countries in each verdict class (supported, partially supported, not supported,
not testable, and the special statuses of H5). The pattern is then labelled:
\begin{itemize}
    \item consistent, when every testable case has the same verdict;
    \item context-dependent, when the testable cases differ;
    \item not assessable, when fewer than two cases are testable.
\end{itemize}
For H1 and H5 the two partial classes are kept distinct, since they point to
different mechanisms. For H4, whose verdict is already comparative, the summary
reports the verdict per technology and the three country profiles on which it
rests. A context-dependent pattern is a result in its own right: it states that
the conclusion does not transfer from one national context to another without
evidence specific to that context.

\subsection{Ranges Instead of Single Values}
\label{subsec:synthesis_ranges}

Every hypothesis statistic is accompanied by its range across the pre-registered
sensitivity variants of its hypothesis (shortlist share, tolerance, climate
configuration, window, cell scale and land scenario where applicable), and by the
number of variants in which the verdict differs from the primary one. The primary
verdict is the one stated; the range shows how far the evidence extends around
it.

% ============================================================================

\section{Scenario Discovery}
\label{sec:prim}

\subsection{Purpose}
\label{subsec:prim_purpose}

Scenario discovery identifies the conditions under which a decision fails, as
regions of the space of uncertain inputs that are simple enough to be
communicated \cite{bryantlempert2010}. In GeoFREA it answers the question that
the hypotheses leave open: under which combinations of costs and climate does the
present-optimal shortlist stop being a good choice? It uses the Patient Rule
Induction Method (PRIM) \cite{friedmanfisher1999}, which searches for a box,
that is, a set of bounds on a few input variables, in which the futures with an
undesirable outcome are concentrated.
% verify friedmanfisher1999

\subsection{Unit of Analysis and Outcome}
\label{subsec:prim_outcome}

The unit of analysis is the future $f=(m,s)$, not the cell. The descriptors of a
future, its parameter values and its climate member, are shared by all cells, so
a cell-by-cell analysis would repeat the same search thousands of times and
return boxes describing individual cells rather than conditions of the system.
The outcome therefore summarises, for each future, what happens to the
present-optimal shortlist as a whole.

For one country and technology, with the nominal shortlist
$S^{\mathrm{nom}}_k$ of Section~\ref{sec:h3}, the tolerance $m=k$, and
$r_{i,f}^{\mathrm{rank}}$ the rank of cell $i$ in future $f$ by relative regret,
the displaced share of the nominal shortlist in future $f$ is
\begin{equation}
    V_f=\frac{\sum_{i\in S^{\mathrm{nom}}_k}q_i\,
    \mathbb{1}\!\left[r^{\mathrm{rank}}_{i,f}>k+m\right]}
    {\sum_{i\in S^{\mathrm{nom}}_k}q_i}.
    \label{eq:prim_outcome}
\end{equation}
A future is a case of interest when
\begin{equation}
    y_f=\mathbb{1}\!\left[V_f\geq\theta\right],\qquad\theta=0.20,
    \label{eq:prim_case}
\end{equation}
that is, when at least one in five present-optimal megawatts falls beyond the
tolerance. The threshold and the tolerance are those of H1b and H3.

The futures are all combinations of the climate members of the core window and
the parameter samples, with the sample size adopted by the convergence protocol
of the techno-economic phase.

\subsection{Descriptors}
\label{subsec:prim_descriptors}

Each future is described by the values of the sampled techno-economic parameters
in its sample, its climate model and emission scenario as categorical variables,
and the capacity-weighted means over $S^{\mathrm{nom}}_k$ of its change factors
of irradiance, temperature and wind. Land parameters are not descriptors, since
land does not vary across futures.

\subsection{Algorithm Settings}
\label{subsec:prim_settings}

The box is found by peeling: at each step, the slice of a descriptor that most
increases the share of cases of interest inside the box is removed. Peeling
removes a fraction $\alpha=0.05$ of the futures in the box at each step, the
customary value. The box must contain at least $5\%$ of all futures (minimum
support), and at most three descriptors may be restricted, so that the box
remains interpretable. Each box of the peeling trajectory is characterised by
\begin{equation}
    \mathrm{density}=\frac{\#\{f\in\mathrm{box}:y_f=1\}}{\#\{f\in\mathrm{box}\}},
    \qquad
    \mathrm{coverage}=\frac{\#\{f\in\mathrm{box}:y_f=1\}}{\#\{f:y_f=1\}}.
    \label{eq:prim_density_coverage}
\end{equation}
Density measures how reliably the box predicts the undesirable outcome; coverage
measures how much of it the box explains.

\subsection{Box Selection Rule}
\label{subsec:prim_selection}

The trajectory trades density against coverage. Choosing a box after inspecting
the trajectory would make the result an adjustment to the data; the selection is
therefore fixed in advance:
\begin{equation}
    \begin{aligned}
        \mathrm{box}^{*}=\arg\max\;\{\,\mathrm{density}:\ & \mathrm{coverage}\geq0.50,\\
        & 1\leq\text{restricted descriptors}\leq3\,\}.
    \end{aligned}
    \label{eq:prim_selection}
\end{equation}
The unrestricted box, which has coverage one by construction, is not eligible.
The selected box explains at least half of the cases of interest and, among the
boxes that do, is the one in which the undesirable outcome is most concentrated.

Special cases are reported as such: if no future is a case of interest, no
vulnerability is found under the futures evaluated; if every future is, the
outcome is not a vulnerability of particular conditions but a general property
of the shortlist; if no box meets the coverage requirement, no simple
explanation is found.

\subsection{Robustness of the Discovered Box}
\label{subsec:prim_robustness}

A box found once is a point estimate. Its stability is assessed by resampling the
futures with replacement $B=200$ times and repeating the search with the same
settings and selection rule. The results report, for the primary box:
\begin{itemize}
    \item the range (5th to 95th percentile over replicates) of each bound, of
    density and of coverage;
    \item the share of replicates in which each descriptor is restricted, which
    shows whether a descriptor is a stable driver or an artefact of the sample.
\end{itemize}
In addition, the search is repeated with the outcome threshold
$\theta\in\{0.10,0.30\}$, the minimum coverage of the selection rule in
$\{0.40,0.60\}$ and the minimum support in $\{0.02,0.10\}$. Boxes are thus reported
as ranges, consistent with the treatment of every other result of the thesis.

\subsection{Limitations}
\label{subsec:prim_limitations}

The box describes the futures evaluated, not the full space of possible futures;
regions outside the parameter ranges and the climate ensemble are not explored.
Boxes are axis-aligned and may describe poorly a vulnerability that depends on a
combination of variables, such as a ratio of two costs. The climate descriptors
are summaries over the shortlist and do not represent the spatial pattern of
change. With six climate models, the climate dimension is coarse, and a climate
model restricted in a box may reflect the behaviour of one model rather than a
physical condition. The PRIM implementation of GeoFREA is checked against a
reference implementation on stored cases and differs from it in documented
details.
